\documentclass[11pt]{article}
\usepackage[margin=1.1in]{geometry}
\usepackage{amsmath,amssymb}
\usepackage{booktabs}
\usepackage{multirow}
\usepackage{natbib}
\usepackage[hidelinks]{hyperref}
\usepackage{xcolor}
\usepackage{microtype}

\title{Persistence Is Not Accumulation:\\ Erasure, Suppression, and Re-instatement\\ in Online In-Weight Memory for Language Models}
\author{Xuhao Lin\\ \small Independent researcher \quad \texttt{linxuhao84@gmail.com}}
\date{July 2026 --- v1 (preprint)}

\newcommand{\qwen}{Qwen3.5-2B}

\begin{document}
\maketitle

\begin{abstract}
Can a frozen language model accumulate new facts \emph{in its weights}, one at a time, as they arrive?
We study single-pass online fact streams written into a LoRA adapter on a frozen \qwen{} base, with a
de-confounded instrument: per-fact recall timelines, a two-alternative recognition probe that separates
\emph{storage} from \emph{expression}, matched compute budgets, and a capability firewall (adapter-off
equals base) verified in every run. Four findings. (1)~\textbf{Accumulation is re-instatement, not
persistence}: median time-to-first-miss is nearly identical (3--8 writes) for every mechanism we test,
while end-of-stream retention varies by $12\times$; what separates mechanisms is the probability that a
missed fact is later recovered (4\%--94\%). Half-life is the wrong abstraction for online memory.
(2)~\textbf{Whether forgetting is erasure or suppression depends on the write protection}: facts
forgotten under naked sequential writes are unrecognizable (2-AFC at chance), whereas facts forgotten
under accumulated elastic-weight consolidation (EWC) remain $\sim$84\% recognizable --- suppressed, not
erased. (3)~\textbf{Error-gated replay converts suppressed facts into expressed ones at zero extra
budget}: spending an unchanged 4-step replay budget on self-test failures first lifts a summed-Fisher
EWC + replay composition from $\sim$71\% to \textbf{85--96\%} of a 48-fact day recallable from weights
alone --- but the same gating barely helps replay without EWC, because there is nothing latent to
re-instate. (4)~\textbf{Online fact writes are readout-proximal}: restricting the adapter to the last
third of the stack retains most capacity, while the first third is nearly useless --- the opposite of the
early-layer placement reported for offline batch knowledge injection, suggesting the answer to ``where
should LoRA memory live?''\ flips between the batch and streaming regimes. All results are 3-seed
directions on a single substrate; we release the instrument and raw per-fact timelines.
\end{abstract}

\section{Introduction}

Language-model agents increasingly ship with external memory --- files, retrieval, context --- while their
weights stay frozen after deployment. A recurring alternative is to let the agent \emph{write} to a
parameter-efficient adapter online, one interaction at a time, and consolidate at rest. Recent
systematic work has audited LoRA-as-memory in the \emph{batch} regime: inject a corpus by multi-epoch
fine-tuning, measure capacity, saturation, and composition \citep{back2026lora}. That audit explicitly
defers the regime that deployed agents actually face: a \emph{time-incremental stream}, one fact per
turn, single pass, no i.i.d.\ shuffling.

This paper is an empirical audit of that streaming regime. The regime is harsh by construction:
sequential SGD has no mechanism for preserving earlier writes, and interference accumulates. But
\emph{how} it is harsh turns out to be structured, and the structure is invisible to the aggregate
accuracy numbers that continual-learning papers usually report. Our instrument therefore records
\emph{per-fact} trajectories: when each fact is first recalled, when it first fails, whether it ever
comes back, and --- separately --- whether it remains \emph{recognizable} when it is no longer
\emph{recallable}.

\paragraph{Contributions.}
\begin{enumerate}
\item \textbf{A de-confounded streaming instrument} (\S\ref{sec:instrument}): 48 collision-free facts
written one per turn into a LoRA on a frozen base; greedy-recall and 2-AFC recognition probes every 2
writes; a pinned recovery statistic; compute-matched mechanism comparisons; and a capability firewall
(adapter-off GSM8K vs.\ base) checked in every run ($+0.00$ in all runs, all mechanisms, all seeds).
\item \textbf{Accumulation is re-instatement} (\S\ref{sec:reinstatement}): across mechanisms spanning
$12\times$ in final retention, median time-to-first-miss varies only from 3 to 8 subsequent writes.
Mechanisms differ almost entirely in recovery-after-miss (4\%--94\%). Survival-time (``half-life'')
framings of online memory measure the wrong thing.
\item \textbf{Erasure vs.\ suppression is mechanism-dependent} (\S\ref{sec:suppression}): recognition
probes show naked-SGD forgetting destroys fact-specific binding (chance-level 2-AFC), while EWC
forgetting leaves $\sim$84\% of facts discriminable behind a failing recall readout. This connects the
machine-unlearning observation that ``forgotten'' knowledge is often latent
\citep{unlearning-deletion,quantization-unlearning,kotha2023} to online continual learning, with a
mechanism-level answer to \emph{when} each case obtains.
\item \textbf{Error-gated replay exploits suppression at zero extra cost} (\S\ref{sec:errorgated}):
with the replay budget held fixed, selecting replay targets by self-test failure lifts the EWC+replay
composition to 41--46/48 (85--96\%) recallable from weights at end of day. The gain is
\emph{conditional on the EWC substrate}; miss-gating plain replay is within noise. We do not claim the
selection rule as novel \citep{aljundi2019mir,sure2025,mssr2026}; the contribution is the budget-matched
decomposition of \emph{why} it works: latent (suppressed) misses are cheap to re-instate; erased misses
are not.
\item \textbf{A batch/streaming regime flip in layer placement} (\S\ref{sec:layers}): at matched
adapter capacity, the last third of the stack retains most of full-stack capacity while the first third
retains almost nothing --- opposite to the early-layer/FFN advantage reported for batch injection
\citep{back2026lora} and to mid-stack concentration reported for RL post-training \citep{onelayer2026}
and locate-then-edit knowledge editing \citep{meng2022rome,meng2023memit}.
\end{enumerate}

We present these as \emph{floors and directions}, not ceilings and points: one substrate, synthetic
facts, 3 seeds per condition on nondeterministic hardware (\S\ref{sec:limits}).

\section{Related work}\label{sec:related}

\textbf{LoRA as knowledge memory (batch regime).} \citet{back2026lora} systematically audit LoRA as a
parametric memory under multi-epoch injection: capacity scales with rank and saturates; QA-format
supervision dominates raw text; multi-LoRA composition is bottlenecked by routing and merging. Their
stated next step --- time-incremental streams where information ``must be stored, retrieved, and revised
continuously'' --- is precisely the regime studied here, and we find that at least one of their
operational conclusions (early-layer placement) inverts in it. \textbf{Continual learning and replay
selection.} Catastrophic forgetting under sequential training is classical
\citep{mccloskey1989,kirkpatrick2017}. Replay selection by estimated interference (MIR;
\citealt{aljundi2019mir}), encoding-time surprise (SuRe; \citealt{sure2025}), forgetting-curve schedules
(MSSR; \citealt{mssr2026}), and prioritized experience replay \citep{schaul2015} all improve on uniform
replay; our miss-gated rule is a member of this family, applied via generative self-testing, and we
claim only its interaction with the substrate. Orthogonal-projection methods \citep{farajtabar2020ogd}
and pseudo-rehearsal \citep{robins1995} bound the design space we sweep. \textbf{Forgetting as
suppression.} In machine unlearning, ``forgotten'' knowledge is frequently recoverable by quantization,
benign relearning, or prompting \citep{quantization-unlearning,unlearning-deletion,kotha2023};
example-level forgetting-and-relearning events are pervasive even in offline training
\citep{toneva2019}. We import the storage/expression distinction \citep{allenzhu2024} into the online
regime with a per-fact recognition meter. \textbf{Memory architectures.} Complementary-systems designs
--- fast weights, sleep-time consolidation, activation-state hippocampi
\citep{behrouz2024titans,ahn2025} --- motivate but are orthogonal to this audit: we characterize the
write dynamics such architectures inherit.

\section{The instrument}\label{sec:instrument}

\paragraph{Substrate.} \qwen{} (24 layers), fp32, frozen. One LoRA adapter (rank 64,
$\alpha{=}128$, all linear modules unless stated) is the sole trainable memory. AdamW, lr $3\times
10^{-5}$, gradient clip 1.0. Hardware: consumer ROCm GPUs; runs are nondeterministic even at fixed
seed, so every number below is reported per-seed (seeds 1234/2025/777) and interpreted as a direction.

\paragraph{Stream.} $N{=}48$ facts of the form \emph{``The user's \{attribute\} is \{value\}.''} with
unique attributes and unique three-syllable pseudoword values (collision-free; unguessable; answers
cannot be produced without the write). One fact per turn, single pass: 8 gradient steps on the fact
(loss on value tokens only), then the mechanism's bookkeeping, then (for replay mechanisms) exactly
$m{=}4$ replay gradient steps on stored past facts. Budgets are identical across mechanisms.

\paragraph{Mechanisms.} \emph{naked}: writes only. \emph{BF}: Benna--Fusi multi-timescale synaptic
cascade \citep{benna2016} applied to adapter parameters. \emph{EWC}: elastic weight consolidation
\citep{kirkpatrick2017} in the accumulating regime --- diagonal Fisher of each fact is \emph{summed}
into a running Fisher (never rebuilt), anchor re-set after each fact, $\lambda{=}300$ (results are
$\lambda$-robust across 100--1000). \emph{replay}: 4 steps/turn on buffer samples, uniform.
\emph{ewc+replay}: both. \emph{Selection policy}: uniform, or \emph{miss} --- replay targets are drawn
first from facts that failed the most recent self-test, filled uniformly. The self-test is the recall
probe itself; the information it uses (which stored facts currently fail) is available to any deployed
agent that keeps a log of what it wrote.

\paragraph{Probes.} Every 2 writes, all facts so far are probed. \emph{Recall}: greedy completion of
the cloze prefix; substring match on the value. \emph{Recognition}: sum log-probability margin of the
true value vs.\ a fixed distractor value sampled from the \emph{same stream} --- the distractor was
written into the same adapter with the same procedure, so a positive margin measures fact-specific
binding, not token familiarity. Chance is 24/48. \emph{Recovery} (pinned definition): a fact is a
recovery candidate if it was recalled at some probe and missed at a later one; it recovered if it is
recalled again after its first miss. \emph{Firewall}: GSM8K (10-item fixed subset) with the adapter
disabled, before and after the stream; any deviation from the base score indicates capability damage.

\section{Results}

\subsection{Accumulation is re-instatement, not persistence}\label{sec:reinstatement}

\begin{table}[t]
\centering\small
\begin{tabular}{lccccc}
\toprule
mechanism & final recall /48 (seeds) & mean & median first-miss & recovery & recognized /48 \\
\midrule
naked            & 4 / 4 / 3    & 3.7  & 3   & 4\%  & 33 / 23 / 26 \\
BF cascade$^\dagger$ & 4 / 3 / 5 & 4.0 & 5   & 10--15\% & --- \\
EWC ($\lambda$300) & 5 / 18 / 9 & 10.7 & 5   & 27\% & 38 / 45 / 38 \\
replay (uniform) & 24 / 16 / 17 & 19.0 & 5   & 76\% & 45 / 40 / 45 \\
replay (miss)    & 15 / 28 / 27 & 23.3 & 4   & 87\% & 45 / 46 / 46 \\
ewc+replay (uniform) & 32 / 42 / 29 & 34.3 & 8 & 76\% & 47 / 48 / 46 \\
ewc+replay (miss) & \textbf{46 / 44 / 41} & \textbf{43.7} & 6 & \textbf{94\%} & \textbf{48 / 48 / 48} \\
\bottomrule
\end{tabular}
\caption{48-fact single-pass stream, probes every 2 writes, 3 seeds per row. Recovery is pooled over
seeds (pinned definition, \S\ref{sec:instrument}); recognition is the final-probe 2-AFC count (chance
$=24$); a strict margin threshold ($>1$ nat) changes no count by more than 4. $^\dagger$BF numbers are
from the same instrument prior to the recognition meter; its rank/timescale sweep is flat (2--5/48)
across $dt\in\{0.1,0.5\}$, cascade depth $\{8,32\}$. Firewall: $+0.00$ in every run.}
\label{tab:main}
\end{table}

Table~\ref{tab:main} spans a $12\times$ range in final retention (3.7 to 43.7 of 48). If mechanisms
differed by making facts \emph{persist} longer, time-to-first-miss would covary with that range. It
does not: the median fact first fails 3--8 writes after acquisition \emph{under every mechanism}. What
covaries with retention is what happens \emph{after} the first miss: under naked writes 4\% of missed
facts ever return; under uniform replay 76\%; under the error-gated composition 94\%. End-of-stream
alive counts (pooled over 3 seeds, /144) tell the same story: naked 11, EWC 32, replay 57, composition
103, error-gated composition 131.

Online in-weight memory in this regime is therefore not a leaky tank whose drain rate mechanisms
adjust; it is a \emph{flickering} store in which nearly everything falls below the recall threshold
quickly, and capacity is determined by the re-instatement process. Aggregate accuracy and half-life
statistics --- the standard reporting in this literature --- cannot see this structure; per-fact timelines
can. The same dynamic is known \emph{inside} offline training as example forgetting events
\citep{toneva2019}: shuffled multi-epoch training is, in effect, a built-in re-instatement schedule.
Our results show the online regime obeys the same law, with the schedule made explicit and paid for.

\subsection{Forgetting is erasure or suppression, depending on the write rule}\label{sec:suppression}

The recognition column of Table~\ref{tab:main} dissociates two kinds of forgetting. Under naked writes,
recognition of the full stream is at or near chance (23--33 vs.\ chance 24): when recall of a fact
dies, its fact-specific binding is gone --- \emph{erasure}. Under accumulated EWC, recall is poor (10.7)
but recognition is far above chance (38--45 of 48, $\sim$84\%): most ``forgotten'' facts remain
discriminable --- \emph{suppression} below the generative readout, not destruction. Replay-based
mechanisms recognize 40--48 while recalling 16--46; the error-gated composition recognizes
\emph{everything} (48/48/48) while recalling 41--46.

This gives a mechanism-level answer to a question the unlearning literature has raised from the other
side --- ``is forgotten knowledge deleted or hidden?''
\citep{unlearning-deletion,quantization-unlearning}: \emph{it depends on what protected the weights
while subsequent writes arrived}. Unprotected interference deletes; Fisher-anchored interference hides.
The practical consequence is \S\ref{sec:errorgated}.

\subsection{Error-gated replay: same budget, $+9$ facts, and why}\label{sec:errorgated}

Holding the replay budget fixed at 4 steps/turn and changing only \emph{selection} --- replay facts that
failed the agent's most recent self-test first --- lifts the composition from 32/42/29 to
\textbf{46/44/41} of 48 (85--96\% of the day held in weights at day's end), with recovery-after-miss
rising from 76\% to 94\%. The same gating applied to replay \emph{without} EWC moves the mean only from
19.0 to 23.3 with overlapping seed ranges and one unstable run.

The recognition data explain the asymmetry. On the EWC substrate, missed facts are suppressed but
latent ($\sim$84\% recognizable), so a targeted gradient step re-instates them cheaply; the miss list
is precisely a list of latent facts. Without EWC, a missed fact is more often erased; replaying it is
re-learning, not re-instatement, and competes for the same interference budget it is trying to escape.
Selection-by-priority is well precedented \citep{schaul2015,aljundi2019mir,sure2025,mssr2026}; the
budget-matched decomposition --- \emph{gating helps exactly insofar as the substrate keeps misses
latent} --- is, to our knowledge, new, and it is the actionable design rule: \textbf{pair a
suppression-inducing protector with an error-gated re-instatement schedule.} In cognitive terms the
composition implements Leitner-style spaced repetition with retrieval practice
\citep{roediger2006} in weight space; the substrate condition is the part the flash-card analogy does
not predict.

\subsection{Online writes are readout-proximal: a regime flip}\label{sec:layers}

\begin{table}[t]
\centering\small
\begin{tabular}{lcccc}
\toprule
& early (0--7) & mid (8--15) & late (16--23) & all 24 \\
\midrule
naked   & 1.0 & 3.3 & 2.7 & 3.7 \\
EWC     & 1.3 \hfill (0\%) & 3.7 \hfill (1\%) & 8.0 \hfill (29\%) & 10.7 \hfill (27\%) \\
replay  & 3.0 \hfill (30\%) & 12.3 \hfill (69\%) & \textbf{21.3} \hfill (74\%) & 19.0--23.5 \hfill (76\%) \\
ewc+replay & --- & 27.7 & 26.0 & 34.3 \\
\bottomrule
\end{tabular}
\caption{Mean final recall /48 (3 seeds; recovery-\% in parentheses) with the adapter restricted to
8-layer bands at matched rank and parameter count (14.5M) vs.\ the full stack ($\sim$43M). The
late-band composition recognizes 48/48/47 while recalling 26 --- the band stores fully and expresses
less.}
\label{tab:layers}
\end{table}

Restricting the adapter to matched-capacity 8-layer bands (Table~\ref{tab:layers}) yields a monotone
depth gradient: the early band is nearly useless for storage under every mechanism; the late band
alone recovers $\sim$90\% of full-stack replay capacity at a third of the parameters, and it is the only
band in which EWC's suppression-and-recovery behavior exists at all (29\% vs.\ 0--1\%). Restricting to
the middle band --- the region favored by locate-then-edit knowledge editing
\citep{meng2022rome,meng2023memit} and by RL layer-contribution analyses \citep{onelayer2026} ---
\emph{hurts} every protected mechanism relative to the full stack.

This is the opposite of the batch-regime placement finding, where early-layer/FFN adapters memorize
best under multi-epoch injection \citep{back2026lora}. The two results are compatible: multi-epoch
optimization can afford to build features early in the stack, while a single-pass gradient on a novel
token binding takes the shortest path to the output margin --- late layers, where our unique pseudoword
targets are nearly orthogonal. We flag the open readout confound: our early/mid band runs predate the
recognition meter, so we cannot yet distinguish ``early bands store nothing'' from ``early bands store
without expressing''; the late-band composition's full recognition at reduced recall shows the meter
can separate these.

\subsection{The firewall holds everywhere}

In every run of every mechanism, band, and policy (including all compositions), disabling the adapter
returns exactly the base model's GSM8K score ($0.70 \to 0.70$ on the fixed 10-item subset). Online
writing at any tested intensity costs the frozen base nothing, by construction --- the architectural
premise (frozen capability, disposable memory cartridges) survives its hardest workload to date.

\section{Discussion}\label{sec:discussion}

\textbf{For memory-agent design.} The numbers argue for a specific division of labor: an external log
(``the file'') as ground truth and \emph{syllabus}, and an in-weight day store trained from it by
error-gated re-instatement --- with the log also supplying the self-tests. The day store is then best
understood not as a transcript but as a \emph{familiarity index}: in the strongest configuration it
recognizes 100\% of the day while recalling 85--96\%, and in weaker configurations the recognition
surplus is exactly what re-grounding against the log can exploit. \textbf{For theory.} A first-order
scale-invariance argument says per-write ``gentleness'' cannot create accumulation (shrinking the step
shrinks a fact's own margin and its interference on others proportionally); only write \emph{geometry}
and \emph{re-instatement schedules} can. Our mechanism ladder is a walk along exactly those two levers,
and linear-attention fast-weight memories --- which meta-learn their write geometry offline
\citep{ahn2025} --- can be read as the activation-state solution to the same problem.
\textbf{For the unlearning debate.} Erasure and suppression are both real and separable by a cheap
2-AFC meter; arguments about whether forgetting ``really'' deletes should condition on the write-time
protection regime.

\section{Limitations}\label{sec:limits}

Single substrate (\qwen{}); results are directions, not calibrated magnitudes, on nondeterministic
ROCm hardware with 3 seeds per condition. The stream is synthetic and collision-free; real facts
correlate, conflict, and get revised. Recall probes are form-matched to the writes (this inflates all
absolute recall numbers; it is held constant across comparisons); paraphrase probing is future work.
The self-test that gates replay is the measurement probe, making probe frequency (here every 2 writes)
a mechanism hyperparameter we have not yet swept. Early/mid band runs lack recognition data. The
firewall subset is small (10 GSM8K items). BF rows predate the recognition meter. The composition's
uniform baseline spans 29--42 across five runs; seed variance is the dominant uncertainty throughout.

\section{Conclusion}

On a single-pass fact stream, in-weight memory does not persist --- it flickers, and it accumulates
exactly to the degree that something re-instates it. Whether the flicker hides or destroys information
is set by the write-time protector, and an error-gated schedule can harvest hidden facts at no extra
cost, to 85--96\% of a day. Where the writes land matters and inverts between batch and streaming
regimes. All of this is invisible to aggregate accuracy; the per-fact instrument is the contribution we
most hope others reuse.

\begin{thebibliography}{99}\small

\bibitem[Aljundi et al.(2019)]{aljundi2019mir} Aljundi, R., et al. Online continual learning with
maximally interfered retrieval. \emph{NeurIPS} 2019.

\bibitem[Allen-Zhu \& Li(2024)]{allenzhu2024} Allen-Zhu, Z., \& Li, Y. Physics of language models:
Part 3.1, knowledge storage and extraction. \emph{ICML} 2024.

\bibitem[Back et al.(2026)]{back2026lora} Back, S., Lee, D., Kang, N., Lee, T., Hong, S.~K., Gwon, Y.,
\& Ahn, S. Understanding LoRA as knowledge memory: An empirical analysis. \emph{ICML} 2026.
arXiv:2603.01097.

\bibitem[Behrouz et al.(2024)]{behrouz2024titans} Behrouz, A., Zhong, P., \& Mirrokni, V. Titans:
Learning to memorize at test time. arXiv:2501.00663.

\bibitem[Benna \& Fusi(2016)]{benna2016} Benna, M.~K., \& Fusi, S. Computational principles of
synaptic memory consolidation. \emph{Nature Neuroscience}, 19(12).

\bibitem[Farajtabar et al.(2020)]{farajtabar2020ogd} Farajtabar, M., et al. Orthogonal gradient
descent for continual learning. \emph{AISTATS} 2020.

\bibitem[Hazard et al.(2025)]{sure2025} Hazard, C., et al. SuRe: Surprise-driven prioritised replay
for continual LLM learning. arXiv:2511.22367.

\bibitem[Hu et al.(2022)]{hu2022lora} Hu, E., et al. LoRA: Low-rank adaptation of large language
models. \emph{ICLR} 2022.

\bibitem[Kirkpatrick et al.(2017)]{kirkpatrick2017} Kirkpatrick, J., et al. Overcoming catastrophic
forgetting in neural networks. \emph{PNAS}, 114(13).

\bibitem[Kotha et al.(2024)]{kotha2023} Kotha, S., Springer, J.~M., \& Raghunathan, A. Understanding
catastrophic forgetting in language models via implicit inference. \emph{ICLR} 2024. arXiv:2309.10105.

\bibitem[McClelland et al.(1995)]{mcclelland1995} McClelland, J.~L., McNaughton, B.~L., \& O'Reilly,
R.~C. Why there are complementary learning systems in the hippocampus and neocortex. \emph{Psych.\
Review}, 102(3).

\bibitem[McCloskey \& Cohen(1989)]{mccloskey1989} McCloskey, M., \& Cohen, N.~J. Catastrophic
interference in connectionist networks. \emph{Psychology of Learning and Motivation}, 24.

\bibitem[Meng et al.(2022)]{meng2022rome} Meng, K., Bau, D., Andonian, A., \& Belinkov, Y. Locating
and editing factual associations in GPT. \emph{NeurIPS} 2022.

\bibitem[Meng et al.(2023)]{meng2023memit} Meng, K., et al. Mass-editing memory in a transformer.
\emph{ICLR} 2023.

\bibitem[MSSR(2026)]{mssr2026} MSSR: Memory-aware adaptive replay for continual LLM fine-tuning.
arXiv:2603.09892.

\bibitem[Robins(1995)]{robins1995} Robins, A. Catastrophic forgetting, rehearsal and pseudorehearsal.
\emph{Connection Science}, 7(2).

\bibitem[Roediger \& Karpicke(2006)]{roediger2006} Roediger, H.~L., \& Karpicke, J.~D. Test-enhanced
learning: Taking memory tests improves long-term retention. \emph{Psychological Science}, 17(3).

\bibitem[Schaul et al.(2015)]{schaul2015} Schaul, T., Quan, J., Antonoglou, I., \& Silver, D.
Prioritized experience replay. \emph{ICLR} 2016. arXiv:1511.05952.

\bibitem[Toneva et al.(2019)]{toneva2019} Toneva, M., et al. An empirical study of example forgetting
during deep neural network learning. \emph{ICLR} 2019.

\bibitem[AHN(2025)]{ahn2025} Artificial hippocampus networks for efficient long-context
modeling. arXiv:2510.07318.

\bibitem[Unlearning(2025)]{unlearning-deletion} Unlearning isn't deletion: Investigating reversibility
of machine unlearning in LLMs. arXiv:2505.16831.

\bibitem[Quantization(2025)]{quantization-unlearning} Catastrophic failure of LLM unlearning via
quantization. \emph{ICLR} 2025. arXiv:2410.16454.

\bibitem[One Layer(2026)]{onelayer2026} Is one layer enough? Training a single transformer layer can match
full-parameter RL training. arXiv:2607.01232.

\end{thebibliography}

\appendix
\section{Reproduction}
Driver: \texttt{phase5/phase5\_accum.py} (mechanism, \texttt{--layers}, \texttt{--replay-policy},
recognition margins logged at every probe); analysis with pinned recovery definition:
\texttt{phase5/analyze\_accum.py}; raw per-fact timelines: \texttt{phase5/results/accum\_*.json}.
Hyperparameters as in \S\ref{sec:instrument}; seeds 1234/2025/777 throughout.

\end{document}
